
\section{The Majority’s Dominion in the Church}

Editorial article in “Folkebladet” for April 20, 1898.

Pastor Kildahl published in “Folkebladet” for March 30 a longer article in defense of the absolute dominion of the majority in the Church, as this is exercised in the United Church through what he calls “the representative system.” It is evident that Kildahl nowhere imagines the possibility of anything other or better than the polity of the United Church; it is this that he defends as the right one, because it is “representative.”\footnote{
From Pastor T. M. Kildahl’s article the following is reproduced:

“… I believe that the old church order is better than the new, because according to the old church order it is the congregations that have joined themselves together into a church body, and it is the congregations that alone, through their representatives, decide all things that concern the church body. I believe further that the old church order is better than the new, because according to the old church order it is the congregation itself that is to own and administer its theological school, choose and depose teachers at the theological school, and in every respect govern it. Thirdly, I believe that the old church order is better than the new, because according to the old church order it is the congregations that hold the right of call, both when it concerns calling theological teachers, and when it concerns calling missionaries. Finally, I believe that the old church order is better than the new, because according to the old church order the right of ordination too rests with the congregations. I believe, therefore, that the old church order is better than the new, because the old church order is wholly throughout lawful and liberal-minded. It safeguards the congregation’s freedom and rights in the affairs of the church body.

I do not believe that the new church order is as good as the old, because according to the new church order the congregations as such are completely deprived of independent standing in all the affairs of the church body. They have nothing to say in the affairs of the church body. They can adopt no resolution concerning the use of the church body’s money. They do not have the right of call when it concerns sending out missionaries for the church body; they cannot choose the church body’s ordaining minister. All these things are decided by people who come together for a meeting, without the congregations having chosen them for that. They do not have the right of call when it concerns choosing teachers at the theological school; they do not own the theological school; they have nothing to say over the theological school. For it all lies in the hands of a board which chooses itself. The congregations have only the right to provide the means, when it concerns the affairs of the church body. …

Therefore in the United Church there is no upper house consisting of pastors and no lower house consisting of laypeople. There is one assembly, where all have the same right, because they all represent congregations. It is the congregations that through their pastors and lay delegates govern and order everything. …

My opinion is this, that in church matters for the congregations there exists no higher authority than the congregations themselves. But I mean that this is the case not only when it concerns the particular local congregation’s own internal affairs, but also when it concerns the matters that the congregations have in common, namely the affairs of the church body. And so it is according to the old church order. The church body is not over the congregations, but the congregations are over the church body; for it is the congregations that govern the church body, but the church body does not govern the congregations. According to the old church order the congregations themselves decide, altogether freely and independently, all matters, both congregational matters and the affairs of the church body. The old church order is therefore free-church, and for that reason I hold to it.

And I am altogether against the new church order as an order dangerous to a free Church and free congregations. For according to the new church order the congregations cannot freely and independently decide all their matters themselves. For according to the new church order the congregations cannot decide the matters that the congregations have in common, that is, the affairs of the church body. They cannot determine how the church body’s money is to be used. They cannot determine who shall teach those who at the theological school are to be educated as pastors for the congregations. They cannot decide how the theological school is to be ordered. They cannot decide whom the church body shall appoint as missionaries to the heathen. They have lost the right of call in all those matters where the congregation in common, that is, as a church body, issues the call. Where the representative system is abolished, there the congregations are deprived of authority in all the affairs of the church body; they have lost their voting right in the affairs of the church body. There are other authorities which are altogether independent of the congregations, and which have not received their authorization from the congregations, which therefore come in one respect to be over the congregations, and which decide all these things. The congregations have only the freedom either to submit to what these authorities determine, or to withdraw. And then it seems to me that a considerable breach has been made in their freedom. One may perhaps object that as many of the members of the congregations as wish may come to the Annual Meeting and exercise influence on the course of things. That is true, but if these are not chosen by the congregations and therefore do not represent the congregations, then it is not the congregations that decide and act. And then the free-church principle is given up.

And then “Folkebladet” says that when I hold the opinion that it is much better to have the government of the church body over the congregations than that the congregations altogether freely and independently decide all their matters themselves, then I surely have the right to hold that opinion. Against that, however, I will protest in the strongest terms. I have often heard this, that a man has the right to hold whatever sort of opinion he pleases, if only he is true to his conviction. But that is a false rationalistic principle. One must indeed often put up with people if they hold opinions that are contrary to the Word of God; but they do not have the right to hold such opinions. I believe that it is contrary to the Word of God that the church body should be over the congregations; therefore I do not believe that anyone, whether “Folkebladet” or I, has the right to hold such an opinion. I believe that it is a high-church, Catholic, and contrary-to-the-Word-of-God opinion that the congregations should not have all authority, right, and freedom also with respect to the government of the church body. And because I believe that the new church order is high-church, Catholic, and contrary to the Word of God, therefore I am against it, and therefore I do not believe that its defenders have the right to hold the opinion they hold.”
}

This system is, in Kildahl’s opinion, the most perfect, because when the congregations have elected, then they are represented, and when they are represented, then the opinion of the representatives is also the opinion of the congregations.

Here we note at once a great omission. God and his Son and his Spirit are altogether omitted from Kildahl’s consideration of church polity. The whole manner of consideration and conception is altogether outward and worldly. It cannot be supposed that the will of God could be different from the will of the majority.

Further, when Kildahl proceeds from the point of view that the majority at the Annual Meetings is the congregations, in this sense namely, that when the majority has determined something, then it is the congregations that have determined it, then that is simply subtlety and sophistry, which Kildahl knows very well. … Kildahl’s assertion, that the determinations of the Annual-Meeting majority are the determinations of the congregations, is altogether untenable, because it is disingenuous.

It is also untenable because it is unreasonable. If, for example, 2 small congregations with 20 eligible voters between them send 4 delegates, and 1 large congregation with 200 eligible voters sends 2 delegates (it can send no more), then those two small congregations can vote for a measure which the large congregation votes against; yet according to Kildahl’s opinion the case is supposed to stand thus, that such a resolution is also the resolution of the large congregation!! Which is nonsense.

Then indeed the Annual Meeting of the United Church is generally a good deal more honest in its statements; it commonly introduces its resolutions with “The church body resolves” or “The Annual Meeting resolves;” one never finds: “The congregations resolve.”

As we said before: God is omitted from this church polity. And therefore conscience too is omitted from this church polity. … We venture to assert that if 100 congregations voted for what was wrong, and one congregation alone for what was right, then the 100 had no right whatever to rule over the conscience of the one.

But suppose now that Kildahl were right, that when someone is chosen to represent the congregation, then the congregation’s understanding, feeling, and will are gathered in him. …

The consequence becomes this, that the most practical thing of all would be for the congregations all together to choose only one representative, or pope, and in him the congregations’ understanding, feeling, and will were then gathered because he was chosen, and when the chairman spoke, then it was the congregations, because he was chosen by them. In this way matters could be decided more easily, more quickly, more cheaply, and more surely, and yet it was always the congregations’ decision; for they themselves had chosen for themselves one great representative!

Yes, “the representative system,” which is a Catholic invention in the Church, has indeed led to the papacy. Well then; if Kildahl can show from Scripture that God has ordained this representative system, then we shall bow and acknowledge that we have been altogether mistaken about the congregations and their right. But be altogether certain of this, that in the congregational question we bow to no other authority.

Georg Sverdrup
